% !PRINTED 498
\weberpage{498}
%<<CONT
from which it follows that for no other value of $u^{24}$ than 64 do two
of the values $v_\nu$ become equal to one another. This equality therefore
occurs only for $\omega = i$ and for certain values of $\omega$ equivalent to $i$,
and hence not for any other purely imaginary value
of $\omega$.

The beginning of the development in ascending powers of
$q$ now gives:
\weq{26}{\begin{aligned}
v_2 - v_5 &= f\left(\frac{\omega - 2\,.\,48}{7}\right) - f\left(\frac{\omega + 2\,.\,48}{7}\right)\\
&= 2\,i\,q^{-\frac{1}{7\,.\,24}} \sin\frac{4\pi}{7} + \cdots,
\end{aligned}}
so that for a sufficiently large real value of $-i\omega$ the
left side of (26) is positive imaginary. If now $-i\omega$ goes
along a real path from $\infty$ to 1, then, as we have seen above,
$v_2 - v_5$ does not pass through zero. Hence for $\omega = i$ also
the difference $v_2 - v_5$ must be positive imaginary, as was assumed
 in (24).

If, then, one substitutes the values (23), (24) into (14), there
results for $\omega = i$:
\[
w_0 = 0, \quad w_5 = 0
\]
\weq{27}{\begin{aligned}
w_1 = w_2 = w_3 = w_6 &= \frac{7 + i\sqrt{7}}{2},\\
w_4 &= -\frac{i\sqrt{7}\,(1+i\sqrt{7})^2}{4}.
\end{aligned}}
Accordingly, for the case $n = 7$ the resolvent can be formed
completely. It reads:
\weq{28}{\begin{aligned}
&w^2\left(w - \frac{7+i\sqrt7}{2}\right)^4\left(w + \frac{i\sqrt7\,(1+i\sqrt7)^2}{4}\right)\\
&= \left(u^{12} - \frac{64}{u^{12}}\right)^2,
\end{aligned}}
and it takes a simpler shape when
\[
w + \frac{i\sqrt7\,(1+i\sqrt7)^2}{4} = z^2
\]
is set:
\weq{29}{z\left(z^2 - \frac{i\sqrt7\,(1+i\sqrt7)^2}{4}\right)\left(z^2+i\sqrt7\right)^2 = u^{12} - \frac{64}{u^{12}}.}
